% ============================================================================
% kapitel/simulation.tex
% ============================================================================

\section{Die Simulation der IWT}

Die IWT wurde numerisch implementiert.

Die Simulation verwendet:

\begin{itemize}
    \item \( N = 4096 \) Knoten
    \item \( 500 \) Iterationen
    \item GPU-Beschleunigung (OpenCL)
    \item Leapfrog-Integrator (symplektisch)
\end{itemize}

\section{Die eingefrorenen Funktionen}

Die Simulation hat zwei eingefrorene Funktionen:

\begin{enumerate}
    \item \textbf{Fluktuationen (Vakuum):} Nur im Vakuum aktiv
    \item \textbf{Energiesenke (Strukturen):} Nur in Strukturen aktiv
\end{enumerate}

Diese Funktionen sind stabil und werden nicht mehr geändert.

\section{Die Parameter}

Die Parameter der Simulation sind:

\begin{align*}
    \text{SCALE} &= 0.70710678 \\
    \text{ALPHA\_0} &= 1e-2 \\
    \text{RHO\_0} &= 1e-6 \\
    \text{RHO\_MIN} &= 1e-8
\end{align*}

\section{Die Stabilität}

Die Simulation ist stabil:

\[
\sum |I|^2 = \text{const}
\]

Axiom 2 ist erfüllt.